\section*{Chapitre \ref{ch:rs:from-signal-to-forecast} --- Du signal à la prévision}

\begin{solution}{exo:rs:from-signal-to-forecast:1}
$0.04 \times 10~\% \times (-1.5) = -0.6~\%$ pour le mois~; $0.04 \times 5~\% \times 2 = 0.4~\%$.
\end{solution}

\begin{solution}{exo:rs:from-signal-to-forecast:2}
$0.03\sqrt{12 \times 500} = 2.32$~; $0.03\sqrt{52 \times 50} = 1.53$, si les prévisions hebdomadaires sont indépendantes, ce qui n'est pas le cas pour celles d'un signal lent.
\end{solution}

\begin{solution}{exo:rs:from-signal-to-forecast:3}
2 et 1 violent l'ordre~: on regroupe en 1.5, 1.5. Les deux 3 sont dans l'ordre~; le 2 suivant le viole~: on regroupe les trois derniers avec lui (2.5), ce qui viole encore le premier 3~:
regroupement des trois en $8/3 = 2.667$. 4 est dans l'ordre. Ajustement~: 1.5, 1.5, 2.667, 2.667, 2.667, 4.
\end{solution}

\begin{solution}{exo:rs:from-signal-to-forecast:4}
$100/(1 + 99 \times 0.05) = 16.8$. Pour la réduire de moitié à 8.4~: $1 + 99\rho = 100/8.4$, $\rho = 0.110$.
\end{solution}

\begin{solution}{exo:rs:from-signal-to-forecast:5}
Qian et Hua~: $0.04/0.05 \times \sqrt{12} = 2.77$. La loi~: $0.04\sqrt{12 \times 704} = 3.68$. La loi suppose que l'IC ne varie que par l'échantillonnage
sur 704 titres indépendants, soit un écart type d'environ $1/\sqrt{704} = 0.038$~; le 0.05 mesuré inclut les mouvements communs des facteurs, que la
loi ignore.
\end{solution}

\begin{solution}{exo:rs:from-signal-to-forecast:6}
Sur une journée, $0.02 \times 2~\% = 0.04~\%$. Sur un mois, la somme des IC journaliers, $0.02 + 0.01 = 0.03$, multipliée par la volatilité journalière~: $0.06~\%$~; et non
21 fois la prévision journalière.
\end{solution}

\begin{solution}{exo:rs:from-signal-to-forecast:7}
\texttt{rs\_forecast.concentrated}~: le \omterm{def:rs:from-signal-to-forecast:breadth}{coefficient de transfert} tombe à 0.59 et le ratio d'information à 2.18, contre 0.89 et 3.25 autour de l'indice de
référence à pondération égale. La plupart des titres ont désormais de très faibles poids dans l'indice, si bien que le portefeuille ne peut les sous-pondérer que de leur faible poids, tandis que les quelques grandes valeurs absorbent les
paris négatifs~: les poids ajustés du risque perdent toute ressemblance avec les prévisions.
\end{solution}

\begin{solution}{exo:rs:from-signal-to-forecast:8}
La largeur compte les prévisions indépendantes~: un signal dont les rangs changent peu d'un jour à l'autre (chapitre~13) donne bien moins de 252 nouvelles prévisions par an,
et 3\,000 titres qui partagent secteurs et facteurs représentent moins de 3\,000 paris indépendants. Le IC journalier d'un signal lent est aussi plus petit que son
IC mensuel, les contraintes du portefeuille abaissent le \omterm{def:rs:from-signal-to-forecast:breadth}{coefficient de transfert}, et trader chaque jour coûte de l'argent. Le 26 de la formule est un plafond que personne n'atteint.
\end{solution}

\begin{solution}{pb:rs:from-signal-to-forecast:1}
\begin{enumerate}
\item $0.05 \times 7.7~\% \times 2 = 0.77~\%$ pour le mois.
\item 0.055.
\item Pente 0.92 et $R^2$ 0.24~\% pour la prévision mise à l'échelle~; pente 0.0039 sur le score brut.
\item Un IC de 0.05 explique environ $0.05^2$ de la variance du rendement de chaque titre~; la prévision reste précieuse parce qu'un portefeuille additionne des centaines de
tels petits edges.
\item Lorsque la relation est visiblement non linéaire dans les queues ou qu'elle sature~; avec assez de données dans les queues pour ajuster des marches de façon fiable.
\item 4.60.
\item 4.92~; l'erreur d'échantillonnage d'un ratio d'information sur dix ans de rendements mensuels est d'environ 0.45, donc oui.
\item Résidus de marché~: moyenne 0.0505, écart type 0.0424. Résidus gaussiens~: écart type 0.0331.
\item 569 prévisions indépendantes par mois, contre 704 titres.
\item 0.86, avec 47~\% des positions ramenées à zéro.
\item 3.54 prédit~; 3.00 réalisé.
\item De 0.98 et 3.91 à 0.25~\% par mois jusqu'à 0.86 et 3.00 à 4~\% (tableau).
\item \textbf{Résultat nommé.} Sur les 4.60 de la loi, le bruit de l'IC sur les résidus de marché (569 titres effectifs sur 704) laisse 4.13, le \omterm{def:rs:from-signal-to-forecast:breadth}{coefficient de
transfert} en position longue seule de 0.86 laisse 3.54, et le bruit des contraintes laisse les 3.00 réalisés~: deux tiers de la promesse.
\item Le réduire à nouveau, proportionnellement à la rotation~: le portefeuille en position longue seule, à une tracking error demandée élevée, est celui qui échange le plus.
\item Le réduire (0.59 dans l'exercice~7)~: de faibles poids dans l'indice laissent peu de marge pour les paris négatifs.
\item Davantage de titres ajoutent des paris corrélés (faible \omterm{def:rs:from-signal-to-forecast:breadth}{largeur effective}), des positions plus petites et moins liquides, et davantage d'erreur d'estimation dans le modèle de risque et
l'IC.
\item La promesse de la loi, la moyenne de l'IC divisée par son écart type, le \omterm{def:rs:from-signal-to-forecast:breadth}{coefficient de transfert} et le ratio d'information réalisé, côte à côte.
\item La recherche améliore l'IC et sa stabilité, ainsi que la \omterm{def:rs:from-signal-to-forecast:breadth}{largeur effective} en trouvant des signaux indépendants~; la construction de portefeuille améliore le \omterm{def:rs:from-signal-to-forecast:breadth}{coefficient de
transfert} et le bruit des contraintes.
\item Elle promet un ratio d'information qui ignore le bruit de l'IC, les contraintes et les coûts, et incite à ajouter des titres ou à trader plus souvent.
\item Le ratio d'information que des prévisions d'une qualité donnée pourraient atteindre si elles étaient indépendantes, sans contraintes et libres de trader.
\end{enumerate}
\end{solution}

\begin{solution}{iq:rs:from-signal-to-forecast:1}
$\mathrm{IR} \approx \mathrm{IC}\sqrt{\mathrm{BR}}$ (Grinold, 1989)~: le ratio d'information du portefeuille optimal croît avec la corrélation des prévisions avec
les rendements réalisés et avec la racine carrée du nombre de prévisions indépendantes par an. Hypothèses~: paris indépendants, IC constant, un
portefeuille moyenne-variance non contraint, absence de coûts.
\end{solution}

\begin{solution}{iq:rs:from-signal-to-forecast:2}
Standardiser le score par date, estimer son IC sur les données passées et poser $\alpha_i = \mathrm{IC}\cdot\omega_i z_i$ avec la volatilité résiduelle du modèle
de risque sur l'horizon~; vérifier la calibration hors échantillon (pente de la régression de Mincer--Zarnowitz proche de 1, \omterm{def:rs:from-signal-to-forecast:curve}{courbe de calibration} par classes).
\end{solution}

\begin{solution}{iq:rs:from-signal-to-forecast:3}
Le bruit de l'IC (son écart type comparé à $1/\sqrt n$~: paris corrélés), le \omterm{def:rs:from-signal-to-forecast:breadth}{coefficient de transfert} des contraintes, le bruit des contraintes, les coûts,
et la question de savoir si l'IC utilisé dans la loi a été estimé dans l'échantillon (in-sample). Mesurer chacun~: la prévision introduite à dessein dans le chapitre a perdu 10~\%, 14~\% et 15~\% aux trois premières étapes.
\end{solution}

\begin{solution}{iq:rs:from-signal-to-forecast:4}
La corrélation entre les poids actifs ajustés du risque du portefeuille et les prévisions ajustées du risque. Les contraintes de position longue seule (surtout face à des indices de référence
concentrés), les limites de rotation, la neutralité sectorielle et factorielle, les limites de position et les plafonds de liquidité la réduisent.
\end{solution}

\begin{solution}{iq:rs:from-signal-to-forecast:5}
Hors échantillon~: régresser les rendements réalisés sur les rendements prévus (ordonnée à l'origine proche de 0, pente proche de 1), tracer la moyenne réalisée en fonction de la moyenne prévue par décile de prévision, comparer par
classe de volatilité et par période, et réajuster la calibration seulement sur les données antérieures à chaque date.
\end{solution}

\begin{solution}{iq:rs:from-signal-to-forecast:6}
Avec $x = w\omega$ normalisé et $u = r/\omega$ de moyenne $\mathrm{IC}\,z$ et bruit indépendant unitaire, le rendement actif a pour moyenne $\mathrm{IC}\,x^\top z =
\mathrm{IC}\sqrt n\,\mathrm{TC}$ et un risque unitaire~; l'annualisation sur des périodes indépendantes donne $\mathrm{TC}\cdot\mathrm{IC}\sqrt{\mathrm{BR}}$. L'indépendance du
bruit échoue en premier~: des facteurs communs font bouger les paris ensemble, ce qui explique que l'IC du chapitre ait eu un écart type plus grand sur les résidus de marché que sur
des résidus gaussiens.
\end{solution}
